\newcommand{\HRule}{\rule{0.95\textwidth}{0.4mm}}

%\begin{titlepage}
\thispagestyle{empty}
\begin{center}
{\bfseries
\MakeUppercase{\Autor} \par
}

\end{center}
\parbox[c][40mm][c]{\linewidth}{
\begin{center}
{\bfseries
\Titulo
}
\end{center}
}

\begin{flushright}
\parbox{0.5\linewidth}{
\linespread{1.0}
\footnotesize \bfseries \Trabalho\ apresentado ao curso de \Curso\ do \Departamento\ da \Universidade\ (\Short) como requisito parcial para obtenção do título de Bacharel em \Curso.
}
\end{flushright}

\vspace{10mm}

\begin{flushright}
\begin{tabular}{l}
     {\bfseries Orientador: \Orientador}  \\
     \\
     {\bfseries Coordenador: \Coordenador}
\end{tabular}
\end{flushright}

\vspace{10mm}

\begin{center}{\bfseries
\begin{flushleft}{ %Para alinhar a esquerda
Nota: \rule{20mm}{0.4mm}

\vspace{5mm}

Aprovado em: \rule{20mm}{0.4mm}/ \rule{20mm}{0.4mm}/ \rule{30mm}{0.4mm}
}\end{flushleft}

\vspace*{\fill}

\HRule\\
\Orientador, \Universidade

\vspace*{\fill}

\HRule\\
\BancaA, \BancaAInst

\vspace*{\fill}

\HRule\\
\BancaB, \BancaBInst

%% Retirar o comentário das próximas 4 linhas caso tenha um quarto membro na banca
% \vspace*{\fill}
%
% \HRule\\
% \BancaC, \BancaCInst
}
\end{center}
%\end{titlepage}